% DRAFT — revised literature section per quality_reports/plans/2026-07-30_litreview-restructure-plan.md
% Not input into main.tex; the live version is literature.tex.
% Decisions taken where the plan said "up to you":
%   Q1: Holzmeister/Candrian/Germann included (S2 first paragraph contrast),
%       Argenton included (S1 credit clause), Steffel kept (S1, one sentence).
%   Q2: Blunden placed in S3 (perception paragraph), Arnestad/Allen in S2 closing beat.
%   Q4: bridging paragraph dropped.
%   Q3 honored: no references beyond the bib; the three keys from the commented
%   control-responsibility block (blunden/arnestad/allen) were verified and added
%   to ProjectAlgorithm.bib on 2026-08-11.
% To adopt: replace the content of literature.tex with this file's content
% (or \input this file instead) and run the full bibtex cycle.

\section{Related Literature} \label{literature}

\subsection*{Responsibility avoidance through delegation}

Delegation has long been studied as an efficiency-enhancing arrangement, with principals delegating to agents who hold better information or specialized skills.\footnote{For an overview of the principal-agent literature see \citet{bolton_contract_2005}.} Alongside these efficiency motives, the literature identifies strategic motives, for instance delegation as a commitment device \citep{schelling_strategy_1960, fershtman_gneezy_2001} or as a way to motivate the agent \citep{aghion_formal_1997}. Our paper builds on a third motive. Principals may delegate because doing so shifts responsibility for the outcome away from themselves, both in their own eyes and in the eyes of others \citep{weaver_politics_1986, fiorina_legislator_1986}. \citet{bartling_shifting_2012} provide the canonical experimental evidence. In their design, a dictator either chooses between a fair and an unfair allocation of an endowment or delegates this choice to another person whose monetary incentives are aligned with the dictator's, and a recipient can assign costly punishment to the dictator and the delegee. Delegation effectively shifts punishment from the dictator to the intermediary, and the prospect of avoiding punishment is a strong motive for the delegation decision itself.

Subsequent work identifies two components of this responsibility shifting. The reduction in the principal's punishment runs through the absence of direct interaction with the affected party rather than through a transfer of blame to a deciding intermediary. \citet{coffman_intermediation_2011} finds that the principal no longer escapes punishment when she interacts directly with the recipient, even when the intermediary remains completely passive, and \citet{oexl_shifting_2013} show that the principal escapes punishment even when delegation necessarily eliminates the fair option. The blame the intermediary absorbs, in turn, attaches to the act of choosing. Randomization devices, which do not choose, absorb less blame than deciding humans and attract less delegation \citep{bartling_shifting_2012, oexl_shifting_2013}.

The motive does not require a punishment stage. \citet{hamman_self-interest_2010} find that principals' transfers to recipients fall sharply under delegation even without any punishment opportunity, indicating that responsibility shifting also operates through felt responsibility and moral wiggle room. The attribution structure extends to the credit side, where delegators receive less credit for fair outcomes, although magnitudes are smaller than for blame \citep{argenton_potters_yang_2023}. Consistent with this asymmetry, \citet{steffel_passing_2016} show that people delegate decisions made on behalf of others more readily than decisions made for themselves, especially when outcomes may be negative. Beyond these settings, the responsibility-shifting motive has been documented in ultimatum bargaining \citep{fershtman_gneezy_2001}, sequential voting \citep{bartling_fischbacher_schudy_2015}, delegated deception \citep{erat_avoiding_2013, gawn_lying_2019, kandul_do_2018}, policy-making \citep{hill_does_2015}, and financial decision-making \citep{freer_friedman_weidenholzer_2024}.

In sum, the canonical literature documents that in settings with self-interested choice, delegation reduces both the principal's punishment and her felt responsibility, in turn motivating delegation. Our setting departs from this canonical design along three dimensions. First, the intermediary is algorithmic rather than human. Second, the outcome falls entirely on a passive recipient rather than partly on the principal herself. Third, the underlying decision involves genuine predictive uncertainty rather than a known fair--unfair choice, removing the transparent self-serving intent that punishment sanctions in the canonical designs. Whether the responsibility-shifting motive retains its force under these conditions is the question we take up.

\subsection*{Delegation to algorithms}

A large empirical literature studies when and why people are willing to use algorithmic decision aids. \citet{dietvorst_algorithm_2015} coin the term \textit{algorithm aversion} for the finding that people prefer their own or other humans' judgment over algorithmic predictions, even after observing that the algorithm performs better.\footnote{Interest in the comparison between human and statistical judgment is long-standing. See \citet{meehl_clinical_1954} and \citet{dawes_robust_1979} for the classic demonstrations that simple statistical rules outperform clinical judgment and are resisted nonetheless.} \citet{logg_algorithm_2019} document the opposite pattern of \textit{algorithm appreciation}. The same split appears in delegation decisions, in which some studies find that people delegate less readily to an algorithm than to a human \citep{dietvorst_algorithm_2015}, while others find a preference for algorithmic over human agents \citep{logg_algorithm_2019, candrian_rise_2022, germann_algorithm_2023, holzmeister_delegation_2023}. Systematic reviews document substantial heterogeneity in algorithm preferences across tasks, contexts, and elicitation methods \citep{burton_systematic_2020, mahmud_what_2022, chugunova_we_2022, qin_ai_2025, irlenbusch_human_2026}.

The determinants documented in this literature attach to three objects. Properties of the algorithm matter. Aversion is stronger when the algorithm cannot be observed to learn from its mistakes \citep{berger_watch_2021} or remains a black box to its users \citep{litterscheidt_financial_2020}, evidence about relative performance gathered through experience and feedback reduces aversion \citep{filiz_reducing_2021}, whereas explanations of how the algorithm works leave acceptance unchanged \citep{dargnies_aversion_2026}. Properties of the task matter. Aversion rises with the degree of irreducible uncertainty \citep{dietvorst_people_2020} and with the perceived subjectivity of the task \citep{castelo_task-dependent_2019}, while reliance on algorithmic advice increases with task difficulty \citep{bogert_humans_2021} and under time pressure \citep{jung_towards_2021}. And the allocation of decision authority matters. People use imperfect algorithms far more readily when they retain even a minimal ability to modify the output \citep{dietvorst_overcoming_2018, normann_delegate_2025}, and part of measured aversion reflects a reluctance to relinquish decision authority as such rather than distrust of algorithms \citep{ivanova_stenzel_delegating_2025, ivanova_stenzel_measuring_2024}.

Authority connects to responsibility. The mere existence of a manual override increases the blame humans receive when the machine errs \citep{arnestad_manual_2024}, and experienced professionals, who answer for outcomes in their domain, are particularly averse to algorithmic advice \citep{allen_algorithm_2022}. Much of this literature concerns advice-taking, in which the decision-maker retains final authority over the outcome. We study the starker margin of full delegation, in which the final prediction is ceded to the algorithm entirely. More importantly, each determinant above attaches to the algorithm, the task, or the decision-maker's authority over the output. Whether the decision-maker is \textit{accountable to the party affected by the decision} is an institutional determinant that this literature has not varied. Preferences for or against algorithms as such produce level effects on delegation. Our between-condition comparison holds the algorithm, the task, and the authority structure fixed and varies accountability alone, so that such level effects difference out.

\subsection*{Responsibility perceptions of decisions involving algorithms}

Within this body of work, a smaller subset focuses on algorithm use in settings whose consequences fall partly or wholly on someone other than the decision-maker, sometimes called the \textit{moral domain} \citep{gogoll_rage_2018}. Much of this work concludes that aversion to algorithms is amplified once third-party welfare is at stake. \citet{gogoll_rage_2018} document this in an incentivized calculation task with performance uncertainty, \citet{bigman_people_2018} in vignettes across high-stakes moral domains, and \citet{jauernig_people_2022} identify \textit{moral discretion} as the human capacity that people specifically value in such settings.\footnote{For recent surveys of this area see \citet{bonnefon_moral_2024} and, with a legal-policy focus, \citet{sunstein_anatomy_2024}.}

Once decisions affect a third party, the question arises how observers apportion responsibility when an algorithm is involved.\footnote{In principle, felt responsibility may be reducible through delegation also in decisions that affect only the decision-maker herself. The evidence discussed here concerns decisions with consequences for others.} A predominantly vignette-based literature finds that human actors are judged more leniently when an algorithm is involved. Physicians who follow a bad computerized recommendation are faulted less than physicians who err unaided \citep{pezzo_pezzo_physician_2006} and face reduced legal exposure when following standard-care AI recommendations \citep{tobia_when_2021}, algorithmic discrimination in hiring triggers less moral outrage than identical human discrimination \citep{bigman_algorithmic_2023}, observers attribute less responsibility to a company when an accident involves technology \citep{kurtzberg_attribution_2004}, and moral violations by algorithms produce less outrage than identical violations by humans across domains \citep{bigman_people_2018}. Attributions are also sensitive to the precise division of authority between human and machine, with humans who monitor algorithms blamed less than humans acting alone \citep{shank_attributions_2019}. Passing a decision on can itself be judged negatively, however, as those asked to decide penalize delegators for transferring decision responsibility \citep{blunden_downside_2023}.

Incentivized evidence is younger, smaller, and divided. One strand finds that algorithmic intermediaries serve as moral cover and attract delegation. AI programmers and delegators escape punishment for inegalitarian AI-mediated allocations, and programmers exploit the resulting moral wiggle room by choosing the inegalitarian option more often \citep{chevrier_algorithm_2024}; delegators of a dictator-game allocation to ChatGPT transfer less to the recipient than non-delegators, indicating strategic use of AI delegation to extract surplus while shedding moral responsibility \citep{tontrup_strategic_2025}; and subjects delegate an incentivized moral donation choice to an AI significantly more often than to a human counterpart \citep{hueholt_trusting_2026}. Another strand finds the opposite pattern, with algorithm aversion dominating blame avoidance in a computer-aided dismissal decision \citep{maasland_blame_2022}. An early comparison finds no significant difference in the dictator's felt responsibility between human-computer and human-human teams \citep{kirchkamp_sharing_2019}. Whether algorithms amplify or dampen the responsibility-shifting motive is thus an open question, and the answer appears to depend on the domain and on how much of the decision is actually transferred.

The most closely related study is \citet{feier_hiding_2022}. In their incentivized experiment, subjects determine a third party's payoff either through their own performance in a Raven-style logic task or, depending on the treatment, by delegating to another human or to an algorithm, while their own payoff is unaffected by this choice. Third parties may then reward or punish the decision-maker. Decision-makers are held \textit{less} responsible for decisions delegated to an algorithm than for decisions delegated to another human, conditional on the outcome, and only algorithmic delegation reduces punishment significantly relative to self-made decisions. Notably, observers treat algorithmic delegation favorably after good outcomes as well, consistent with a general approval of the delegation decision rather than with blame shielding alone.

Our design extends \citet{feier_hiding_2022} along three dimensions. First, their algorithm is calibrated at the session level, so that the delegation decision reflects beliefs about own ability relative to that distribution and thereby carries a signal about the decision-maker's self-perception that may contaminate downstream punishment as a measure of responsibility shifting per se. We calibrate the algorithm at the individual level, removing this relative-performance signal by construction.\footnote{Indeed, their own analysis shows that the propensity to delegate is driven primarily by self-assessed errors in the logic task, not by the type of delegate. Punishment of a non-delegator may then partly sanction perceived overconfidence rather than responsibility-taking per se.} Second, the same subject acts as both delegator and evaluator in their design, potentially introducing self-justification, hypocrisy aversion, and similar motives into the punishment decision. We separate the delegator and evaluator roles across disjoint subject pools.\footnote{For example, in their data, the evaluator's own delegation decision is significantly correlated with the evaluation gap between delegated and non-delegated decisions after good outcomes ($r=-0.34$ in the machine treatment; own computations from their published data).} Third, their design does not include a treatment in which punishment is ruled out by construction, so that the motivational question of whether the prospect of punishment itself drives delegation is not identified. We add this manipulation directly, isolating the role of anticipated punishment in the delegation decision.

Our experiment is thus, to our knowledge, the first to provide a clean test of blame shifting onto an algorithm in a setting with inherent uncertainty, controlling for relative-performance concerns. It is also the first to vary the punishment possibility for an algorithmic delegate, and hence to test whether the prospect of punishment itself motivates the delegation decision.
